% !PRINTED 399
\weberpage{399}
From this it follows:

1. \gesperrt{For a negative $D$ there are in general only the two units}
\[ \varepsilon = +1, \qquad \varepsilon = -1, \]
\gesperrt{and for two special cases}
\[
\begin{aligned}
D=-4\colon\quad &\varepsilon=+1,-1,+i,-i,\\
D=-3\colon\quad &\varepsilon=+1,-1,\\
&-\frac{1+\sqrt{-3}}{2},-\frac{1-\sqrt{-3}}{2},+\frac{1+\sqrt{-3}}{2},+\frac{1-\sqrt{-3}}{2}
\end{aligned}
\]

The numbers of the order $O$ different from zero are then associated in twos, or, in the two exceptional cases, in fours or in sixes.

2. \gesperrt{If $D$ is positive, then there is a solution of the Pell equation:}
\[ T^2 - U^2 D = \pm 4, \]
\gesperrt{in which $T$ and $U$ have positive values that are as small as possible. If then}
\[ \varepsilon = \frac{T + U\sqrt{D}}{2}, \]
\gesperrt{all units are contained in the form $\pm \varepsilon^{\pm\nu}$, where $\nu$ runs through the sequence of numbers $0, 1, 2, \ldots$.}

If equation (2) is solvable for the negative sign, then $N(\varepsilon) = -1$, and $\pm \varepsilon^{\pm\nu}$ has a positive norm for even $\nu$ and a negative norm for odd $\nu$. But if (2) is solvable only for the upper sign, then all units have positive norm. Among the numbers $y$ of the order $O$ which are not units, however, there are always such with negative norm.

If $N(\varepsilon) = -1$, then for every number $y$ in $O$ with negative norm there is an associated one with positive norm. But if $N(\varepsilon) = +1$, then the numbers $y$ split into two classes, of which the one has positive and the other negative norm, and no number of the one class is associated with one of the other class.

Accordingly we consider only the numbers $y$ of the order $O$ with positive norm, among which $y$ and $\pm \varepsilon^{\pm\nu} y$ are associated only when
\[ \varepsilon = \frac{T + U\sqrt{D}}{2} \]
%CONT>>
